
\section{Fixed traces, loops, and quantifier-free invariants}
\label{sec:trace}

\subsection{State coding}

A finite program state is encoded by a natural number.  Tuples are nested
through the checked pairing function; bounded arrays may use a standard
primitive-recursive sequence code.  The representation exports projections
and proofs that packing followed by projection returns each component.

For a theorem-specific loop, fix:
\[
   s_0(\vec p),\qquad
   \mathsf{step}(t,s,\vec p),\qquad
   \mathsf{run}(t,\vec p),
\]
where \(\mathsf{run}\) is defined by primitive recursion:
\[
\begin{aligned}
 \mathsf{run}(0,\vec p)&=s_0(\vec p),\\
 \mathsf{run}(t+1,\vec p)
   &=\mathsf{step}(t,\mathsf{run}(t,\vec p),\vec p).
\end{aligned}
\]
The current \code{recDef}, \code{rec_zero}, and \code{rec_succ} APIs construct
the symbol and its two defining equations from open terms.

\subsection{The generic invariant theorem}

Let \(\Inv(s,\vec p)\) be quantifier free.  Define
\[
      J(t,\vec p)\defeq
       \Inv(\mathsf{run}(t,\vec p),\vec p).
\]
Substitution of terms into a quantifier-free formula preserves
quantifier-freeness, so \(J\) is an admissible induction formula.

\begin{theorem}[Fixed trace reflection]
Assume object-iPRA derivations of
\[
\begin{aligned}
 &\Inv(s_0(\vec p),\vec p),\\
 &\Inv(s,\vec p)
    \ \imp\
   \Inv(\mathsf{step}(t,s,\vec p),\vec p).
\end{aligned}
\]
Then iPRA derives
\[
       \Inv(\mathsf{run}(T,\vec p),\vec p)
\]
for arbitrary \(T\).
\end{theorem}

\begin{proof}
Use quantifier-free induction on \(J(t,\vec p)\).

At zero, rewrite \(\mathsf{run}(0,\vec p)\) to \(s_0(\vec p)\) using
\code{rec_zero}, then apply the base derivation and equality substitution.

For the successor step, assume
\(\Inv(\mathsf{run}(t,\vec p),\vec p)\).  Specialize the local step derivation
at
\[
     s=\mathsf{run}(t,\vec p).
\]
This gives
\[
 \Inv(\mathsf{step}(t,\mathsf{run}(t,\vec p),\vec p),\vec p).
\]
Rewrite the state to \(\mathsf{run}(t+1,\vec p)\) using
\code{rec_succ}.  The result is \(J(t+1,\vec p)\).

Apply \code{IPRA.ind} or \code{PRA.qfInd} and instantiate the universally
derived conclusion at \(T\).
\end{proof}

\begin{remark}
The theorem is parameterized at the Lean metalevel by the concrete terms
\(s_0\), \(\mathsf{step}\), and \(\Inv\).  It is not an object-language
universal interpreter.  The generated primitive-recursive symbol is specific
to the loop.
\end{remark}

\subsection{Terminal extraction}

Often the invariant alone is not the desired theorem.  Add a terminal
contract
\[
  \Inv(s,\vec p),\ \mathsf{halt}(s,\vec p)=1
  \ \vdash_{\PRAAx}\
  \Post(\mathsf{out}(s,\vec p),\vec p).
\]
If a clock theorem proves
\[
   \mathsf{halt}(\mathsf{run}(T,\vec p),\vec p)=1,
\]
then simultaneous cut with fixed trace reflection gives the postcondition.

The clock may be:
\begin{itemize}
\item a literal input bound;
\item a rank or lexicographic measure with explicit primitive-recursive fuel;
\item a search capacity;
\item a precision schedule computed by a backward demand pass;
\item a residual producer obligation.
\end{itemize}
Strict descent by itself is not enough when no primitive-recursive envelope is
known.  The clock is a first-class proof boundary.

\subsection{Early stopping without partiality}

A bounded loop that should stop when a condition becomes true is made total by
freezing its state:
\[
 \mathsf{step}'(t,s)=
 \begin{cases}
   s,&\mathsf{halt}(s)=1,\\
   \mathsf{step}(t,s),&\mathsf{halt}(s)=0.
 \end{cases}
\]
The branch is a primitive-recursive selector built from the Boolean halt bit.
The invariant includes stability of a halted state.  The run is total for
every clock, while a separate theorem proves that the requested clock is long
enough.  This avoids partial functions inside the object language.

\subsection{Changing parameters and backward demands}

Some loops consume a parameter that changes with time: precision \(k_t\),
remaining budget \(b_t\), or a finite support demanded by future nodes.  Encode
all changing parameters in the state.  The invariant is then a formula about
the complete packed state.

For a finite acyclic computation, demands can be computed backward.  If node
\(v\) requires input precision
\[
  d_{u\leftarrow v}(q_v)
\]
from predecessor \(u\), define
\[
  \Need_u=\max_{v:u\to v} d_{u\leftarrow v}(\Need_v).
\]
The maximum over the finite outgoing list is primitive recursive.  A backward
run computes all \(\Need_v\); a forward run performs the actual calculations
at those fixed demands.  The backward run has its own quantifier-free
invariant:
every processed edge has been assigned enough precision for its consumer.

For a cycle, a demand update need not stabilize.  PCFM permits two safe forms:
\begin{enumerate}
\item an explicit cap and a proof that the capped value is sufficient;
\item a primitive-recursive iteration bound and a post-fixed-point
certificate.
\end{enumerate}
It does not infer a fixed point from monotonicity alone.

\subsection{Nested loops and flattening}

Nested bounded loops may be handled compositionally or flattened.

\paragraph{Compositional method.}
Prove the inner loop as a module parameterized by the outer state.  Use that
module in the outer step proof.  This is often easiest to read.

\paragraph{Flattened method.}
Pack the outer and inner program counters and registers into one state.  Define
one transition and one clock.  This can reduce repeated binder and
specialization work in Lean.

Both methods yield the same kind of final derivation.  There is no requirement
that every theorem use a single universal normal form.  PCFM treats
normalization as an engineering choice.

\subsection{Ranked finite descent}

Let states satisfy a Boolean validity condition and let
\(\rho(s)\) be a natural rank.  Suppose a nonterminal transition strictly
decreases rank and all states remain within a primitive-recursive envelope.
Set the clock to the initial rank plus one.

\begin{theorem}[Ranked termination]
If
\[
  \neg\mathsf{halt}(s)\imp
  \rho(\mathsf{step}(s))<\rho(s),
\]
then the frozen run is halted by time \(\rho(s_0)+1\).
\end{theorem}

\begin{proof}
Use the invariant
\[
  \mathsf{halt}(\mathsf{run}(t))=1
  \ \vee\
  \rho(\mathsf{run}(t))+t\le \rho(s_0).
\]
All relations are quantifier free.  At an unhalted step the rank decreases by
at least one, so the sum with time does not increase.  At time
\(\rho(s_0)+1\), the second disjunct would imply a natural number is strictly
below zero; therefore the run is halted.
\end{proof}

Lexicographic descent is reduced to a single rank only after component bounds
are known.  For bounded coordinates \(a_i<B_i\), mixed-radix encoding gives
\[
  \rho(a_0,\ldots,a_{r-1})
    = a_0+B_0(a_1+B_1(\cdots)).
\]
If the number of coordinates or their bounds are not fixed by primitive-
recursive data, no fuel theorem is claimed.

\subsection{Trace certificates versus recomputation}

A producer may supply a complete trace \(z_0,\ldots,z_T\).  A verifier checks
the initial state, each local transition, and the terminal property.  This is
often preferable when recomputing a large calculation in the object term
would be expensive.  The trace checker is a bounded Boolean fold, and its
reflection theorem gives:
\[
  \CertCheck(\vec z)=1
   \ \vdash_{\PRAAx}\
  z_T=\mathsf{run}(T).
\]
The certificate does not create trust: local transition equality is checked
inside iPRA.  It changes only the cost profile and allows external search to be
separated from internal verification.
